\section{A Frobenius--Koszul obstruction to the formal data}
\label{frob:section}
\label{sec:frobenius}

The obstruction in this section is an algebraic statement about formal
power series.  Its hypotheses include both ideal-membership laws: neither
law follows here merely from the vanishing of the potential on a graph.
The proof will take place in coordinate Frobenius quotients, with the
ordinary differential on forms as well as the wedge Koszul differential.

\begin{theorem}[Nonexistence of the formal data]
\label{frob:no-data}
Let $\Bbbk$ be a field of prime characteristic $\ell>2$, and let
$d\geq5$.  For $1\leq i\leq d$, choose an integer $n_i\geq0$, a tuple
$a_i=(a_{i,1},a_{i,2},a_{i,3})$ of variables, and a tuple $y_i$ of
$n_i$ variables.
Put $D=\sum_i n_i$, write $a=(a_1,\ldots,a_d)$, and let
$C=(C_i^+,C_i^-)_{i=1}^d$ be $2D$ variables, with $n_i$ variables in
each indicated block.  There do not exist series
\[
 P_i\in(a_i,y_i)^3\Bbbk[[a_i,y_i]],\qquad
 Y=(Y_i^+,Y_i^-)_{i=1}^d\in
       (a,t)\Bbbk[[a,t]]^{2D},\qquad t=(t_1,\ldots,t_D),
\]
having the following properties.  Define
\begin{equation}
 \Psi(a,C)=\sum_{i=1}^d
       \bigl(P_i(a_i,C_i^+)-P_i(0,C_i^-)\bigr),\qquad
 p_i(y_i)=(\partial_{a_{i,1}}P_i)(0,y_i),
 \label{frob:separated}
\end{equation}
and let $\sigma$ interchange the $+$ and $-$ blocks for each $i$.
All indicated derivatives are taken before substitution.
First, the two linear maps
\[
 A=\partial_tY(0,0),\qquad \sigma A:
                  \Bbbk^D\longrightarrow\Bbbk^{2D}
\]
have complementary images and are injective, or equivalently the square
matrix $[A\ \sigma A]$ is invertible.  Second,
\begin{equation}
                         \Psi(a,Y(a,t))=0.
 \label{frob:graph-zero}
\end{equation}
Finally, for every $i,j,k$ (including repetitions), both memberships
\begin{equation}
 \begin{split}
 a_{i,1}a_{j,1}a_{k,1}&\in J_Y,\\
 p_i(Y_i^-(a,t))p_j(Y_j^-(a,t))p_k(Y_k^-(a,t))&\in J_Y,
 \qquad
 J_Y=\bigl((\partial_{C_\nu}\Psi)(a,Y(a,t))\bigr)_{\nu=1}^{2D}
 \end{split}
 \label{frob:memberships}
\end{equation}
hold in $\Bbbk[[a,t]]$.
\end{theorem}

Only distinct triples will be used in the proof.  The formulation with
all triples records the stronger input that is available. The field
need not be perfect.

The proof has four stages. First the two complementary zero graphs
give differential forms with nonzero residue pairing. The cubic
law for the $p_i$ then selects a tensor class with three suitable positive
blocks. Deforming those blocks over three square-zero parameters lifts
that tensor to an actual cycle. Finally the cubic law for the $a_{i,1}$ makes
the parameter product times the moving graph's normal form a boundary.
The resulting pairing forces a nonzero parameter product to
vanish. The following subsections construct each form and primitive
in the same finite coordinate complex.

\subsection{Forms, normal cycles, and residue}
\label{frob:forms-subsection}

Fix an ordering $C_1,\ldots,C_{2D}$, in block order
$1^+,1^-,2^+,2^-,\ldots,d^+,d^-$.  Define
\begin{equation}
 T=\Bbbk[[C]]/(C_1^\ell,\ldots,C_{2D}^\ell),\qquad
 \Omega^q=T\otimes_{\Bbbk}
                 \bigwedge^q\langle\mathrm dC_1,\ldots,\mathrm dC_{2D}\rangle.
 \label{frob:quotient}
\end{equation}
The coefficient monomials $C^e$ with $0\leq e_\nu<\ell$ form a basis
of $T$.  The ordinary formal partial derivatives descend to $T$, since
$\partial_{C_\mu}(C_\nu^\ell)=0$.  Consequently the usual differential
$\mathrm d$ is well-defined on $\Omega$, satisfies $\mathrm d^2=0$, and
obeys the graded product rule.  These are also the ordinary K\"ahler
forms of $T$ over $\Bbbk$: the relations $C_\nu^\ell=0$ impose no
relations on the coordinate differentials.

For $F\in\Bbbk[[C]]$ put
\[
 \delta=\mathrm dF\wedge{},\qquad H=H^*(\Omega,\delta).
\]
We have $\delta^2=0$ and
$\mathrm d\delta+\delta\mathrm d=0$.  Thus ordinary differentiation
induces a degree-one differential $\mathcal D$ on $H$:
$\mathcal D[\omega]=[\mathrm d\omega]$.  This differential must be
distinguished from $\delta$, whose cohomology defines $H$.

The finite coordinate calculation has a classical antecedent in the
coordinate proof of Cartier's theorem: the one-variable classes
$C_\nu^{\ell-1}\mathrm dC_\nu$ and their tensor products survive
ordinary differentiation
\cite[Theorem~7.2 and (7.2.11)--(7.2.12)]{Katz1970}.
Here the additional differential $\delta=\mathrm dF\wedge{}$ and the
two graph membership laws will connect those forms to the obstruction.
Their required properties are proved below.

Define $\operatorname{Res}_C$ on $\Omega$ by taking, in top form degree,
the coefficient of $C_1^{\ell-1}\cdots C_{2D}^{\ell-1}$ in the
coefficient of $\mathrm dC_1\wedge\cdots\wedge\mathrm dC_{2D}$; it is
zero in other degrees.  Directly from the monomial basis,
\begin{equation}
                         \operatorname{Res}_C(\mathrm d\omega)=0.
 \label{frob:residue-exact}
\end{equation}
Indeed a derivative producing exponent $\ell-1$ in its differentiated
variable would have to differentiate exponent $\ell$, whose coefficient
is zero in characteristic $\ell$; in the chosen basis no such monomial
is present.  If $\eta$ is homogeneous and $\delta\zeta=0$, then
\begin{equation}
 (\delta\eta)\wedge\zeta
   =(-1)^{|\eta|}\eta\wedge\delta\zeta=0.
 \label{frob:boundary-wedge}
\end{equation}
It follows that
\[
 \langle[\omega],[\zeta]\rangle
                   =\operatorname{Res}_C(\omega\wedge\zeta)
\]
is a well-defined pairing on Koszul cohomology.  We require no assertion
that this scalar is unchanged by a change of coordinates.

\begin{lemma}[Normal cycles and restricted memberships]
\label{frob:normal-cycle}
Suppose that $\Gamma$ is a formal smooth $D$-dimensional graph in the
$C$ variables, centered at the origin, with $F|_\Gamma=0$.  Complete a
parametrization to formal coordinates $(t,u)$, with $\Gamma$ given by
$u_1=\cdots=u_D=0$.  In $\Omega$ its normal form
\begin{equation}
 \gamma=\left(\prod_{b=1}^D u_b^{\ell-1}\right)
                    \mathrm du_1\wedge\cdots\wedge\mathrm du_D
 \label{frob:normal-form}
\end{equation}
is annihilated by the normal ideal $(u)$ and satisfies
$\mathrm d\gamma=\delta\gamma=0$.
If an ambient function $g$ has a restricted expression
\begin{equation}
  g|_\Gamma=\sum_{\nu=1}^{2D}h_\nu(t)
                               (\partial_{C_\nu}F)|_\Gamma,
 \label{frob:restricted-membership}
\end{equation}
then $g\gamma$ is a Koszul boundary.  More precisely, lift $h_\nu$ to
$\widetilde h_\nu(C)$ by the inverse tangential coordinates and let
$\iota_\nu$ be contraction with $\partial/\partial C_\nu$.  Then
\begin{equation}
 g\gamma=\delta\left(\sum_\nu
                         \widetilde h_\nu\iota_\nu\gamma\right).
 \label{frob:primitive-general}
\end{equation}
If two such graphs have complementary tangent spaces, their normal
forms have nonzero residue pairing.
\end{lemma}

\begin{proof}
A centered formal coordinate change and its inverse preserve the
coordinate Frobenius ideals.  For a centered series
$f(z)=\sum_{e\ne0}c_ez^e$ in characteristic $\ell$, the coefficientwise
Frobenius identity gives
\[
 f(z)^\ell=\sum_{e\ne0}c_e^\ell z^{\ell e}
                          \in(z_1^\ell,\ldots,z_{2D}^\ell).
\]
Apply this to every component of each inverse coordinate map.  Hence
we may calculate in $\Bbbk[[t,u]]/(t^\ell,u^\ell)$, where the
notation means the separate $\ell$th powers of all coordinates.
Multiplication by $u_b$ annihilates \eqref{frob:normal-form}.  Its ordinary
differential is zero because every term differentiating its coefficient
repeats one of its normal differentials.  Write $F(t,u)\in(u)$.  The
tangential coefficients of $\mathrm dF$ belong to $(u)$, whereas each
normal differential wedges to zero with the full normal volume.
Therefore $\delta\gamma=0$.

The difference
$g-\sum_\nu\widetilde h_\nu\partial_{C_\nu}F$ lies in $(u)$ before
truncation and continues to do so after truncation.  The contraction
identity in the original coordinates is the coefficient-linear
extension of the exterior and interior product identity in
\cite[Section~2, p.~72, (2.9)]{Koszul1950}:
\begin{equation}
 \delta\iota_\nu+\iota_\nu\delta
                         =(\partial_{C_\nu}F)\operatorname{id}.
 \label{frob:contraction}
\end{equation}
Since $\delta$ is linear over $T$, annihilation of $(u)$ and
$\delta\gamma=0$ prove \eqref{frob:primitive-general}.  There are no
derivatives of the coefficients $\widetilde h_\nu$ in this calculation.

For the last assertion let $u$ and $v$ be the respective normal
coordinates.  Complementarity of the tangent spaces says exactly that
$(u,v)$ has invertible linear part, so it too is a full coordinate
system.  The socle of $T$, namely the annihilator of its maximal ideal,
is the one-dimensional space spanned by
$C_1^{\ell-1}\cdots C_{2D}^{\ell-1}$.  This follows directly by
multiplying a monomial expansion by each $C_\nu$.  The coordinate
automorphism carries its nonzero socle generator to
\[
             q=\prod_b u_b^{\ell-1}\prod_b v_b^{\ell-1}\ne0.
\]
The wedge of the two normal forms is $qJ(C)$ times the original
coordinate volume, where $J(C)$ is the unit Jacobian determinant of
$(u,v)$.  Since $q$ is in the socle, $qJ(C)=J(0)q\ne0$.  Its residue
is therefore nonzero.  This proves nonvanishing in the fixed original
coordinates without a coordinate-invariance formula for the residue.
\end{proof}

\subsection{A chain retraction and a tensor-kernel calculation}
\label{frob:linear-subsection}

The next two lemmas address the obstruction to deforming a Koszul
cycle. For a single block, let $Q$ and $p$ be scalar functions in its
Frobenius quotient, and write $\delta=\mathrm dQ\wedge{}$.
Replacing $Q$ by $Q+sp$ over $\Bbbk[s]/(s^2)$ changes the differential
to $\delta+s\,\mathrm dp\wedge{}$. A $\delta$-cycle $\omega$ lifts to
a cycle $\omega+s\xi$ precisely when
\[
                    \delta\xi=-\mathrm dp\wedge\omega.
\]
On the block Koszul cohomology let $S$ be multiplication by $p$ and
let $\mathcal D$ be induced by ordinary differentiation. The product
rule gives $[\mathcal D,S]=[\mathrm dp\wedge{}]$. For a class already
in $N=\ker S$, the lift condition is therefore
$S\mathcal D[\omega]=0$: the classes we need belong to
$K=N\cap\mathcal D^{-1}N$. The cubic law will first express a class
as a sum of tensors with enough factors in $N$. The lemmas below
replace those factors by ones in $K$ through a homotopy for
$\mathcal D$, preserving the pairing with an ordinary-closed normal
cycle.

\begin{lemma}[Retraction with an explicit homotopy]
\label{frob:retraction}
Let $(V,\mathcal D)$ be a finite-dimensional graded complex over a
field, and let $S:V\to V$ have degree zero.  No relation between $S$
and $\mathcal D$ is assumed.  Set
\[
 N=\ker S,\qquad K=N\cap\mathcal D^{-1}N.
\]
There are a degree-zero chain idempotent $r$ and a degree-minus-one
map $h$ such that
\begin{equation}
 r(N)=K,\qquad r|_K=\operatorname{id},\qquad
                  1-r=\mathcal Dh+h\mathcal D.
 \label{frob:retraction-identity}
\end{equation}
The image of $r$ is not asserted to equal $K$.
\end{lemma}

\begin{proof}
The subspace $K$ is a subcomplex: if $x\in K$, then
$\mathcal Dx\in N$ and $\mathcal D^2x=0\in N$.  Choose a graded
complement $N=K\oplus E$.  If $e\in E$ and $\mathcal De\in N$, then
$e\in K$, hence $e=0$.  Thus $\mathcal D$ is injective on $E$, and
$\mathcal DE\cap N=0$.  In particular $K\oplus E\oplus\mathcal DE$
is a graded direct sum.  Also $E\oplus\mathcal DE$ injects into $V/K$:
if $e+\mathcal De'\in K$, then $\mathcal De'\in N$, so $e'=e=0$.

Take any graded vector-space complement $C'$ to that direct sum, and
define
\[
 h(\mathcal De)=e\quad(e\in E),\qquad
                     h|_{K\oplus E\oplus C'}=0.
\]
Set $Q=\mathcal Dh+h\mathcal D$.  The image of $\mathcal Dh$ lies in
$\mathcal DE$, and the image of $h\mathcal D$ lies in $E$, even when
$\mathcal DC'$ has components in the other summands.  Moreover $Q$
is the identity on $E$ and on $\mathcal DE$, and is zero on $K$.
Consequently $Q^2=Q$.  The equation $\mathcal D^2=0$ gives
\[
 \mathcal DQ=\mathcal Dh\mathcal D=Q\mathcal D.
\]
Now $r=1-Q$ has all the asserted properties.  This calculation is why
an arbitrary complement is allowed; an arbitrary vector-space
projection by itself would not give a chain map.
\end{proof}

\begin{lemma}[Joint kernels in a tensor product]
\label{frob:tensor-kernel}
For finite-dimensional vector spaces $V_i$ and endomorphisms $S_i$,
put $N_i=\ker S_i$, and let $W$ be an additional tensor factor on which
all $S_i$ act as the identity.  On
$V_1\otimes\cdots\otimes V_d\otimes W$ one has
\begin{equation}
 \bigcap_{|I|=3}\ker\left(\prod_{i\in I}S_i\right)
 =\sum_{|J|\geq d-2}
   \left(\bigotimes_{i=1}^d V_i(J)\right)\otimes W,
 \qquad
 V_i(J)=\begin{cases}N_i&i\in J,\\ V_i&i\notin J.\end{cases}
 \label{frob:tensor-kernel-identity}
\end{equation}
This also holds with additional passive factors interspersed in any
fixed tensor order.
\end{lemma}

\begin{proof}
Choose complements $V_i=N_i\oplus M_i$.  The total tensor product is
the direct sum of subspaces $V_A$, indexed by subsets
$A\subset\{1,\ldots,d\}$, using $M_i$ when $i\in A$ and $N_i$ when
$i\notin A$, together with the passive factor.  Fix a triple $I$.
The operator $S_I=\prod_{i\in I}S_i$ kills
$\bigoplus_{I\not\subset A}V_A$.  Its restriction to the complementary
subspace, which has $M_i$ in every slot $i\in I$ and unrestricted
factors elsewhere, is injective: it is a tensor product of the
injections $S_i|_{M_i}$ and identity maps over a field.  Thus
\[
                  \ker S_I=\bigoplus_{I\not\subset A}V_A.
\]
Intersecting these coordinate direct sums over all triples leaves
exactly $|A|\leq2$, which is the right-hand side of
\eqref{frob:tensor-kernel-identity}.  The images $S_i(M_i)$ may intersect
$N_i$; this does not affect injectivity on the specified complementary
\emph{domain} subspace.  Passive factors have no effect on the argument.
\end{proof}

We shall also use the following explicit tensor homotopy.  Suppose
each graded complex $(V_b,\mathcal D_b)$ has $r_b,h_b$ as in
\eqref{frob:retraction-identity}; identity maps and zero homotopies are
allowed.  With the usual signed tensor differential, put
\begin{equation}
 R=\bigotimes_b r_b,\qquad
 \mathcal H=\sum_b
      \left(\bigotimes_{c<b}r_c\right)\otimes h_b\otimes
      \left(\bigotimes_{c>b}1\right).
 \label{frob:tensor-homotopy}
\end{equation}
On a homogeneous tensor the $b$th summand in $\mathcal H$ has sign
$(-1)^{\sum_{c<b}|x_c|}$.  Since every $r_c$ is a degree-zero chain
map, terms in which the differential acts in a slot other than $b$
cancel in $\mathcal D\mathcal H+\mathcal H\mathcal D$.  The remaining
terms give the telescoping identity
\begin{equation}
 \mathcal D\mathcal H+\mathcal H\mathcal D
 =\sum_b r_{<b}\otimes(1-r_b)\otimes1_{>b}=1-R.
 \label{frob:tensor-homotopy-identity}
\end{equation}

\subsection{The special fiber and selection of a tensor}
\label{frob:special-fiber}

\begin{proof}[Proof of Theorem~\ref{frob:no-data}]
Suppose the stated data exist.  If $D=0$, the ideal $J_Y$ is zero, so
the first membership in \eqref{frob:memberships} already contradicts
the nonzero monomial $a_{1,1}a_{2,1}a_{3,1}$.  We may assume $D>0$.
Set
\[
 Q_i(y_i)=P_i(0,y_i),\qquad
 F(C)=\Psi(0,C)=\sum_i\bigl(Q_i(C_i^+)-Q_i(C_i^-)\bigr),
\]
and use the complexes and residue in \eqref{frob:quotient}.
The two graphs
\[
               \Gamma_+=Y(0,t),\qquad \Gamma_-=\sigma Y(0,t)
\]
have complementary tangent spaces.  They are smooth formal graphs
because their tangent maps are injective and can be completed to
invertible coordinate maps.  The function $F$ vanishes on both: this
follows from \eqref{frob:graph-zero} and $F(\sigma C)=-F(C)$.

Choose once and for all a constant matrix $L$ whose columns complement
the image of $A$, and use the inverse coordinates of
\begin{equation}
                       C=Y(0,t)+Lu
 \label{frob:fixed-complement}
\end{equation}
to define the normal cycle $\alpha$ for $\Gamma_+$.  Choose any fixed
ordered normal coordinates for $\Gamma_-$ and define its normal cycle
$\beta$.  Lemma~\ref{frob:normal-cycle} gives
\begin{equation}
 \delta\alpha=\delta\beta=\mathrm d\alpha=\mathrm d\beta=0,
 \qquad \langle[\alpha],[\beta]\rangle\ne0.
 \label{frob:initial-pair}
\end{equation}

The second membership in \eqref{frob:memberships}, at $a=0$, transfers
to $\Gamma_-$ as follows.  Regard $\sigma$ also as the permutation of
scalar-coordinate indices.  Differentiating $F(\sigma C)=-F(C)$ gives
\[
 (\partial_{C_\nu}F)(Y(0,t))
       =-(\partial_{C_{\sigma(\nu)}}F)(\sigma Y(0,t)).
\]
On $\Gamma_-$, the positive block $C_i^+$ equals $Y_i^-(0,t)$.
Thus permuting the derivative entries and negating their witnesses
turns that membership into
\[
 \bigl(p_i(C_i^+)p_j(C_j^+)p_k(C_k^+)\bigr)|_{\Gamma_-}
               \in\bigl((\partial_{C_\nu}F)|_{\Gamma_-}\bigr)_\nu.
\]
The explicit primitive in Lemma~\ref{frob:normal-cycle} therefore proves
\begin{equation}
        p_i(C_i^+)p_j(C_j^+)p_k(C_k^+)[\beta]=0
                    \quad\hbox{in }H
 \label{frob:triple-annihilation}
\end{equation}
for every distinct triple.  This is an ambient Koszul boundary
statement, not just an identity after restriction to the graph.

The separated expression for $F$ identifies $(\Omega,\delta)$ with
the graded tensor product of the $2d$ block complexes.  A positive
block has differential $\delta_i^+=\mathrm dQ_i\wedge{}$ and a negative
block has differential $\delta_i^-=-\mathrm dQ_i\wedge{}$.  Under the
identification a tensor is sent to the wedge of its factors in the
fixed block order; both the Koszul differential and ordinary
differentiation carry the sign $(-1)^{\sum_{c<b}|x_c|}$ in block $b$.
Consequently there is a natural identification
\begin{equation}
 H=\bigotimes_{i=1}^d(H_i^+\otimes H_i^-),\qquad
 \mathcal D=\hbox{the signed tensor differential induced by ordinary }
                 \mathrm d_i^\pm.
 \label{frob:kunnet}
\end{equation}
For completeness, the cohomology assertion uses only field linear
algebra: split each finite complex into representatives for its
cohomology with zero differential and pairs $e,\delta e$.  Each pair
is contractible, and its tensor product with any complex is
contractible by the signed tensor homotopy.  The remaining tensor of
the representative spaces is precisely the displayed cohomology.
Naturality of the tensor cycle map gives the asserted ordinary
differential on it.

On $H_i^+$ let $S_i$ be multiplication by $p_i$, and let $\mathcal D_i$
be the differential induced by ordinary differentiation.  Define
\begin{equation}
 N_i=\ker S_i,\qquad K_i=N_i\cap\mathcal D_i^{-1}N_i.
 \label{frob:ki}
\end{equation}
Multiplication by $p_i$ commutes with the block Koszul differential.
The degree-one operator $\mathrm dp_i\wedge{}$ anticommutes with that
differential and induces on $H_i^+$ the commutator
\begin{equation}
 [\mathcal D_i,S_i]=\mathcal D_iS_i-S_i\mathcal D_i
                           =[\mathrm dp_i\wedge{}].
 \label{frob:commutator}
\end{equation}
This operator vanishes on $K_i$, since both $S_i x$ and
$S_i\mathcal D_i x$ vanish there.

Let $z=[\beta]$.  By \eqref{frob:triple-annihilation} and
Lemma~\ref{frob:tensor-kernel}, $z$ lies in the sum of tensor subspaces
having at least $d-2$ positive slots in $N_i$.  Apply the retractions
of Lemma~\ref{frob:retraction} in all positive slots, taking $r_b=1$
and $h_b=0$ in negative slots.  Their tensor $R$ sends that sum into
the analogous sum with at least $d-2$ slots in $K_i$.
Since $\mathcal Dz=0$, \eqref{frob:tensor-homotopy-identity} gives
\[
                           Rz-z=-\mathcal D(\mathcal H z).
\]
Pairing with $[\alpha]$ kills this ordinary-differential boundary.
Indeed if $\omega$ is any Koszul-cycle representative of a class,
then $\mathrm d\alpha=0$ and \eqref{frob:residue-exact} imply
\begin{equation}
 \operatorname{Res}_C(\alpha\wedge\mathrm d\omega)
       =(-1)^D\operatorname{Res}_C\bigl(\mathrm d(\alpha\wedge\omega)\bigr)
       =0.
 \label{frob:pair-homotopy}
\end{equation}
Changing representatives adds a Koszul boundary, which pairs to zero
by \eqref{frob:boundary-wedge}; ordinary differentiation of a changed
representative adds another Koszul boundary by anticommutation.
Thus \eqref{frob:pair-homotopy} is genuinely an identity on $H$.
It follows from \eqref{frob:initial-pair} that
$\langle[\alpha],Rz\rangle\ne0$.

All subspaces and complements used above are graded, and $Rz$ has
form degree $D$.  Express it as a finite sum of homogeneous simple
tensors in the indicated sum of subspaces, keeping only total degree
$D$.  One such tensor, denoted $k$, satisfies
\begin{equation}
             \langle[\alpha],k\rangle\ne0
 \label{frob:selected-pair}
\end{equation}
and has at least $d-2\geq3$ positive factors in their $K_i$.
Fix three of these labels $i_1,i_2,i_3$.  The selected tensor need
not be $\mathcal D$-closed.  What will be used is that each of its
homogeneous factors has a homogeneous cycle representative for its
\emph{block Koszul differential}.

\subsection{Three parameters and all mixed terms}
\label{frob:three-parameters}

Set
\begin{equation}
 B=\Bbbk[s_1,s_2,s_3]/(s_1^2,s_2^2,s_3^2),\qquad
 \mathfrak m_B=(s_1,s_2,s_3),\qquad w=s_1s_2s_3.
 \label{frob:base}
\end{equation}
Substitute $a_{i_r,1}=s_r$ for $r=1,2,3$, and set all other slope
coordinates to zero.  Write this substitution as $a=a(s)$.
By the separated form \eqref{frob:separated} and $s_r^2=0$, exactly
\begin{equation}
 W_s(C):=\Psi(a(s),C)
       =F(C)+\sum_{r=1}^3s_r p_{i_r}(C_{i_r}^+).
 \label{frob:deformed-potential}
\end{equation}
There are no mixed-parameter terms in this potential: each summand
$P_i$ involves only its own $a_i$.  Products of distinct parameters
do survive in $B$ and must be retained in a cycle lift.

From now on the common ambient complex is
\begin{equation}
 T_B=B[[C]]/(C_\nu^\ell)_\nu,
 \qquad \Omega_B=T_B\otimes_B\bigwedge_B\langle\mathrm dC_\nu\rangle,
 \qquad \delta_s=\mathrm d_CW_s\wedge{}.
 \label{frob:common-complex}
\end{equation}
Every differential and contraction here is relative to $B$; in
particular $\mathrm d_Cs_r=0$.  Choose homogeneous block-cycle
representatives $x_b$ for every factor of $k$.  In each selected
positive slot $i_r$, \eqref{frob:commutator} and membership in $K_{i_r}$
give a homogeneous form $y_{i_r}$ of the same degree as $x_{i_r}$ with
\begin{equation}
 \delta_{i_r}^+y_{i_r}=-\mathrm dp_{i_r}\wedge x_{i_r}.
 \label{frob:factor-correction}
\end{equation}
Thus
\[
 (\delta_{i_r}^++s_r\mathrm dp_{i_r}\wedge{})
                         (x_{i_r}+s_ry_{i_r})=0.
\]
Use $x_{i_r}+s_ry_{i_r}$ in the selected slots and $x_b$ in every
other slot, and take their ordered tensor product.  This defines
$\kappa_s\in\Omega_B^D$.  To display every term, suppress the unchanged
factors while retaining them in their original positions in the actual
tensor, and abbreviate the three selected pairs by $x_r,y_r$:
\begin{equation}
\begin{split}
 \kappa_s={}&x_1\otimes x_2\otimes x_3
 +s_1y_1\otimes x_2\otimes x_3
 +s_2x_1\otimes y_2\otimes x_3
 +s_3x_1\otimes x_2\otimes y_3\\
 &+s_1s_2y_1\otimes y_2\otimes x_3
 +s_1s_3y_1\otimes x_2\otimes y_3
 +s_2s_3x_1\otimes y_2\otimes y_3
 +s_1s_2s_3y_1\otimes y_2\otimes y_3.
\end{split}
 \label{frob:eight-terms}
\end{equation}
No permutation of the actual factors is being made in this display.
If $\kappa_A$ denotes the coefficient of $s_A=\prod_{r\in A}s_r$,
the coefficient of $s_A$ in $\delta_s\kappa_s$ is
\begin{equation}
 \delta\kappa_A+
      \sum_{r\in A}(\mathrm dp_{i_r}\wedge{})_{i_r}
                                     \kappa_{A\setminus\{r\}},
 \label{frob:lift-coefficients}
\end{equation}
where slot operators carry the usual preceding-degree sign.  In its
first term only the differentials acting on a $y_{i_r}$ can be
nonzero; by \eqref{frob:factor-correction} these cancel exactly the
corresponding terms in the sum.  Their signs agree because $y_{i_r}$
and $x_{i_r}$ have the same form degree.  An additional perturbation
in a slot already carrying $s_r$ is zero by $s_r^2=0$.  This proves
the cancellation for every subset $A$, including $|A|=2$ and $|A|=3$.
We have therefore constructed an actual cycle with the required
specified reduction:
\begin{equation}
 \delta_s\kappa_s=0,\qquad
         \kappa_s\bmod\mathfrak m_B=\kappa_0,
 \qquad [\kappa_0]=k.
 \label{frob:lifted-cycle}
\end{equation}

\subsection{The moving graph in the same coordinate quotient}
\label{frob:moving-graph}

Keep the exact matrix $L$ and ordering used to define $\alpha$ in
\eqref{frob:fixed-complement}.  Consider
\begin{equation}
                \Theta_s(t,u)=Y(a(s),t)+Lu.
 \label{frob:moving-coordinate-map}
\end{equation}
This has a formal inverse over $B$ despite its possible nonzero
nilpotent constant term in the $C$ variables.  To see this without
assuming that it is centered in those variables, first regard
$s_1,s_2,s_3$ as unrestricted formal variables over $\Bbbk$.
The augmented map
\[
              (s,t,u)\longmapsto (s,Y(a(s),t)+Lu)
\]
is centered at the full origin and has invertible Jacobian: its
diagonal blocks are the identity in $s$ and $[A\ L]$.  The formal
inverse function construction, solving each homogeneous degree using
this invertible linear matrix, supplies an inverse fixing $s$.
It descends modulo $(s_1^2,s_2^2,s_3^2)$.  The resulting maps are
continuous in the $(\mathfrak m_B,C)$-adic and
$(\mathfrak m_B,t,u)$-adic topologies.  Since $\mathfrak m_B^4=0$,
these topologies are equivalent to the respective coordinate-adic
topologies.  Write the inverse coordinates as
\[
                         (t,u)=(T_s(C),U_s(C)).
\]
Their $C$-constant terms belong to $\mathfrak m_B$, because the
reduced inverse at $s=0$ is centered.

Every $b\in\mathfrak m_B$ has $b^\ell=0$.  Indeed write it as a
$\Bbbk$-linear combination of nonempty squarefree monomials in
the $s_r$ and apply Frobenius; the $\ell$th power of each monomial
contains a factor $s_r^2$.  Thus for any component of either the
forward or the inverse coordinate map, written as
$f(z)=b+\sum_{e\ne0}c_ez^e$, we have
\begin{equation}
 f(z)^\ell=b^\ell+\sum_{e\ne0}c_e^\ell z^{\ell e}
                         \in(z_1^\ell,\ldots,z_{2D}^\ell).
 \label{frob:both-directions}
\end{equation}
The two already inverse formal maps therefore induce inverse
isomorphisms
\begin{equation}
 B[[C]]/(C_\nu^\ell)_\nu
       \simeq B[[t,u]]/(t_1^\ell,\ldots,t_D^\ell,u_1^\ell,\ldots,u_D^\ell).
 \label{frob:translated-quotient}
\end{equation}
This argument checks the constant shifts and every nonlinear term
in both directions.  It uses the actual base $B$, not an assertion
about arbitrary nilpotent rings.  Relative forms and their ordinary
differentials transport across the isomorphism; all the displayed
Frobenius relations have zero relative differential.

Define the moving normal form in the original coordinates by
\begin{equation}
 \alpha_s=\left(\prod_{b=1}^D U_{s,b}(C)^{\ell-1}\right)
             \mathrm d_CU_{s,1}(C)\wedge\cdots\wedge\mathrm d_CU_{s,D}(C)
                       \in\Omega_B^D.
 \label{frob:moving-normal}
\end{equation}
The graph identity \eqref{frob:graph-zero} says
$W_s(\Theta_s(t,0))=0$, so $W_s\circ\Theta_s\in(u)$.
The calculation of Lemma~\ref{frob:normal-cycle}, now over $B$, gives
\begin{equation}
 (U_s)\alpha_s=0,\qquad
            \mathrm d_C\alpha_s=\delta_s\alpha_s=0.
 \label{frob:moving-normal-identities}
\end{equation}
Inversion of the augmented formal coordinate map commutes with
reduction modulo $\mathfrak m_B$, by uniqueness of the inverse.
Relative differentiation also commutes with reduction.  Because the
same $L$ and the same normal-coordinate order were retained, this
gives the precise equality
\begin{equation}
                       \alpha_s\bmod\mathfrak m_B=\alpha.
 \label{frob:alpha-reduction}
\end{equation}
It is the originally chosen form, with its original Jacobian factors.

\subsection{An ambient primitive and a unit residue}
\label{frob:contradiction}

Use the first membership in \eqref{frob:memberships} for the selected
triple and substitute $a=a(s)$.  There are $h_\nu(s,t)\in B[[t]]$ with
\begin{equation}
 w=\sum_{\nu=1}^{2D}h_\nu(s,t)
                   (\partial_{C_\nu}W_s)(Y(a(s),t)).
 \label{frob:w-restricted}
\end{equation}
Lift the witnesses to the ambient formal ring by
$\widetilde h_\nu(s,C)=h_\nu(s,T_s(C))$.  Before Frobenius truncation,
subtraction of the $u=0$ value in adapted coordinates gives
\begin{equation}
 w-\sum_\nu\widetilde h_\nu(s,C)\partial_{C_\nu}W_s(C)
                                      \in(U_{s,1},\ldots,U_{s,D}).
 \label{frob:w-discrepancy}
\end{equation}
It remains true after truncation.  In the adapted quotient the kernel
of evaluation $u=0$ is exactly $(u)$, as the monomial basis shows.
That evaluation is legitimate in the original quotient by
\eqref{frob:translated-quotient}; formal partial differentiation is
legitimate there because derivatives of $C_\nu^\ell$ vanish.

Let $\iota_\nu$ now be the original $C$-coordinate contraction on
$\Omega_B$.  The relative version of \eqref{frob:contraction} reads
\[
 \delta_s\iota_\nu+\iota_\nu\delta_s
                         =(\partial_{C_\nu}W_s)\operatorname{id}.
\]
Consequently the explicit ambient form
\begin{equation}
 \eta_s=\sum_{\nu=1}^{2D}\widetilde h_\nu\iota_\nu\alpha_s
                  \in\Omega_B^{D-1}
 \quad\hbox{satisfies}\quad
                          \delta_s\eta_s=w\alpha_s.
 \label{frob:w-primitive}
\end{equation}
Here \eqref{frob:w-discrepancy} is killed by
\eqref{frob:moving-normal-identities}.  Again no derivatives of the
witnesses appear, because $\delta_s$ is coefficient-linear wedge
multiplication.  This boundary identity holds in precisely the same
complex \eqref{frob:common-complex} as \eqref{frob:lifted-cycle}.

Extend $\operatorname{Res}_C$ to $\Omega_B$ by the identical original
coordinate-coefficient rule, now $B$-linearly.  Its reduction modulo
$\mathfrak m_B$ is the previous residue.  By
\eqref{frob:alpha-reduction}, \eqref{frob:lifted-cycle}, and
\eqref{frob:selected-pair}, the scalar
\[
                 b_s=\operatorname{Res}_C(\alpha_s\wedge\kappa_s)
\]
has reduction
\begin{equation}
 b_0=\operatorname{Res}_C(\alpha\wedge\kappa_0)
                         =\langle[\alpha],k\rangle\ne0.
 \label{frob:unit-reduction}
\end{equation}
Hence $b_s$ is a unit in $B$.  Explicitly, if $b_s=b_0+n$ with
$n\in\mathfrak m_B$, then
\[
                b_s^{-1}=b_0^{-1}\sum_{q=0}^3(-n/b_0)^q.
\]
On the other hand the exact wedge identity
\[
 (\delta_s\eta_s)\wedge\kappa_s
              =(-1)^{D-1}\eta_s\wedge\delta_s\kappa_s=0
\]
and \eqref{frob:w-primitive} give
\[
                  0=\operatorname{Res}_C((\delta_s\eta_s)\wedge\kappa_s)
                    =w b_s.
\]
Multiplication by $b_s^{-1}$ forces $w=0$, but the eight squarefree
monomials in $s_1,s_2,s_3$ form a $\Bbbk$-basis of $B$, and
$w=s_1s_2s_3$ is one of them.  This is the contradiction.

\end{proof}

The contradiction takes place entirely in the finite coefficient ring
$B$ and the complex $\Omega_B$. In particular, no flatness of the
deformed Koszul cohomology is required. The order-three hypothesis on
the potentials records the input of Section~\ref{sec:algebra}; the
proof above uses only separation, the zero-graph identity,
complementary tangents, and the two cubic membership laws.
